\section*{根系}

\begin{enumerate}
\def\labelenumi{\arabic{enumi})}
\setcounter{enumi}{4}
\item
  对 \(\alpha\)（\(\mathfrak{h}^*\)——\(\mathfrak{h}\) 的对偶——中的任一元素），令 \(\mathfrak{g}^{\alpha}\) 为由 \(x\in \mathfrak{g}\)（满足 \([h,x] = \alpha (h)x\) 对所有 \(h\in \mathfrak{h}\)）所成的集合。若 \(\alpha = 0\)，则 \(\mathfrak{g}^{\alpha} = \mathfrak{h}\)。任何 \(\alpha \in \mathfrak{h}^{*} - \{0\}\)（若它使 \(\mathfrak{g}^{\alpha}\neq 0\)）都称为 \((\mathfrak{g},\mathfrak{h})\) 的根。以 \(\mathrm{R}(\mathfrak{g},\mathfrak{h})\)（或简记为 R）表示 \((\mathfrak{g},\mathfrak{h})\) 的根的集合。它是 \(\mathfrak{h}^*\) 中第VI章§1第4节意义下的既约根系。\(\mathfrak{g}\) 单当且仅当 R 不可约。
\item
  对所有 \(\alpha \in \mathbb{R}\)，\(\mathfrak{g}^{\alpha}\) 的维数为 1。向量空间 \([\mathfrak{g}^{\alpha}, \mathfrak{g}^{-\alpha}]\) 包含于 \(\mathfrak{h}\)，维数为 1，且含有唯一的元素 \(H_{\alpha}\)（满足 \(\alpha(H_{\alpha}) = 2\)）；我们有 \(H_{\alpha} = \alpha^{\vee}\)（第VI章§1第1节）；诸 \(H_{\alpha}\)（对 \(\alpha \in \mathbb{R}\)）所成的集合是逆根系 \(\mathbb{R}^{\vee}\)（\(\mathbb{R}\) 的）。
\item
  我们有 \(\mathfrak{g} = \mathfrak{h} \oplus \bigoplus_{\alpha \in \mathrm{R}} \mathfrak{g}^{\alpha}\)。存在一个族 \((X_{\alpha})_{\alpha \in \mathrm{R}}\)，使得对所有 \(\alpha \in \mathrm{R}\)，有 \(X_{\alpha} \in \mathfrak{g}^{\alpha}\) 且 \([X_{\alpha}, X_{-\alpha}] = -H_{\alpha}\)。每个 \(x \in \mathfrak{g}\) 都可唯一地写成形式
\end{enumerate}

\[
x = h + \sum_ {\alpha \in \mathrm{R}} \lambda_ {\alpha} X _ {\alpha}, \quad \text { where } h \in \mathfrak {h}, \lambda_ {\alpha} \in k.
\]

两个元素的括号可用下列公式计算

\[
\begin{array}{c} [ h, X _ {\alpha} ] = \alpha (h) X _ {\alpha} \\ \relax [ X _ {\alpha}, X _ {\beta} ] = 0 \quad \text {if} \alpha + \beta \notin \mathrm{R} \cup \{0 \} \\ \relax [ X _ {\alpha}, X _ {- \alpha} ] = - H _ {\alpha} \\ \relax [ X _ {\alpha}, X _ {\beta} ] = \mathrm{N} _ {\alpha \beta} X _ {\alpha + \beta} \quad \text {if} \alpha + \beta \in \mathrm{R}, \end{array}
\]

其中 \(\mathrm{N}_{\alpha \beta}\) 是 \(k\) 中的非零元素。

\begin{enumerate}
\def\labelenumi{\arabic{enumi})}
\setcounter{enumi}{7}
\tightlist
\item
  令 B 为 R 的基。代数 g 由诸 \(X_{\alpha}\) 与诸 \(X_{-\alpha}\)（对 \(\alpha \in B\)）生成。我们有 \([X_{\alpha}, X_{-\beta}] = 0\)（若 \(\alpha, \beta \in B\) 且 \(\alpha \neq \beta\)）。令 \((n(\alpha, \beta))_{\alpha, \beta \in B}\) 为 R 的（关于 B 的）嘉当矩阵。我们有 \(n(\alpha, \beta) = \alpha(H_{\beta})\)。若 \(\alpha, \beta \in B\) 且 \(\alpha \neq \beta\)，则 \(n(\alpha, \beta)\) 是负整数，且
\end{enumerate}

\[
(\operatorname{ad} X _ {\beta}) ^ {1 - n (\alpha , \beta)} X _ {\alpha} = 0 \quad \text { and } \quad (\operatorname{ad} X _ {- \beta}) ^ {1 - n (\alpha , \beta)} X _ {- \alpha} = 0.
\]

\begin{enumerate}
\def\labelenumi{\arabic{enumi})}
\setcounter{enumi}{8}
\tightlist
\item
  若 \(\alpha, \beta, \alpha + \beta \in \mathbb{R}\)，令 \(q_{\alpha \beta}\) 为最大整数 \(j\)（使 \(\beta - j\alpha \in \mathbb{R}\)）。7) 中的族 \((X_{\alpha})_{\alpha \in \mathbb{R}}\) 可以取得使 \(\mathrm{N}_{\alpha, \beta} = \mathrm{N}_{-\alpha, -\beta}\)（当 \(\alpha, \beta, \alpha + \beta \in \mathbb{R}\) 时）。此时 \(\mathrm{N}_{\alpha \beta} = \pm (q_{\alpha \beta} + 1)\)。存在一个对合自同构 \(\theta\)（\(\mathfrak{g}\) 的），它把 \(X_{\alpha}\) 变为 \(X_{-\alpha}\)（对所有 \(\alpha \in \mathbb{R}\)）；我们有 \(\theta(h) = -h\)（对所有 \(h \in \mathfrak{h}\)）。\(\mathbf{Z}\)-子模 \(\mathfrak{g}_{\mathbf{Z}}\)（\(\mathfrak{g}\) 中由诸 \(H_{\alpha}\) 与诸 \(X_{\alpha}\) 生成的）是 \(\mathbf{Z}\)-李子代数（\(\mathfrak{g}\) 的），且典则映射 \(\mathfrak{g}_{\mathbf{Z}} \otimes_{\mathbf{Z}} k \to \mathfrak{g}\) 是同构。
\end{enumerate}

偶对 \((\mathfrak{g},\mathfrak{h})\) 可由一个分裂半单 Q-李代数经标量扩张得到。

\begin{enumerate}
\def\labelenumi{\arabic{enumi})}
\setcounter{enumi}{9}
\item
  R 的 Weyl 群、权群、\ldots{}（以及 \ldots{} 等诸群）称为 \((\mathfrak{g},\mathfrak{h})\) 的相应各群。以下把 Weyl 群记作 W。我们考虑它不仅在 \(h^{*}\) 上、而且（通过结构的转移）也在 h 上的作用。若以 \(h_{Q}\)（分别地，\(h_{Q}^{*}\)）表示 h（分别地，\(h^{*}\)）中由诸 \(H_{\alpha}\)（分别地，诸 \(\alpha\)）生成的 Q-向量子空间，则 h（分别地，\(h^{*}\)）可典则地与 \(h_{Q} \otimes_{Q} k\)（分别地，\(h_{Q}^{*} \otimes_{Q} k\)）等同起来，且 \(h_{Q}^{*}\) 可与 \(h_{Q}\) 的对偶等同起来。当我们谈到 R 的 Weyl 房时，理解为它们在 \(h_{R} = h_{Q} \otimes_{Q} R\) 或 \(h_{R}^{*} = h_{Q}^{*} \otimes_{Q} R\) 中。
\item
  令 \(\varPhi\) 为 \(\mathfrak{g}\) 的 Killing 型。若 \(\alpha + \beta \neq 0\)，则 \(\mathfrak{g}^{\alpha}\) 与 \(\mathfrak{g}^{\beta}\) 关于 \(\varPhi\) 正交。\(\varPhi\) 在 \(\mathfrak{g}^{\alpha} \times \mathfrak{g}^{-\alpha}\) 上的限制是非退化的。若 \(x, y \in \mathfrak{h}\)，则 \(\varPhi(x, y) = \sum_{\alpha \in \mathrm{R}} \alpha(x)\alpha(y)\)。我们有 \(\varPhi(H_{\alpha}, H_{\beta}) \in \mathbf{Z}\)。\(\varPhi\) 在 \(\mathfrak{h}\) 上的限制非退化且在 W 下不变；它在 \(\mathfrak{h}_{\mathbf{Q}}\) 上的限制是正定的。
\item
  \((\mathfrak{g},\mathfrak{h})\) 的根系在同构的意义下只依赖于 g 而不依赖于 h。滥用术语，把 Weyl 群、权群、\ldots{}（\((\mathfrak{g},\mathfrak{h})\) 的）也称为 g 的 Weyl 群、权群、\ldots{}。
\end{enumerate}

若 \(R_{1}\) 是既约根系，则存在分裂半单李代数 \((\mathfrak{g}_{1},\mathfrak{h}_{1})\) 使 \(\mathrm{R}(\mathfrak{g}_{1},\mathfrak{h}_{1})\) 同构于 \(R_{1}\)；且它在同构意义下唯一。

于是可分裂半单李代数的分类就化为根系的分类。

